\subsection{A fixed successful candidate suffices at the literal cutoff}
A further simplification concerns enumeration completeness, not the machine.
The established compiler supplies a finite word
\[
 U=\mathrm{STORE}_0;\mathrm{STORE}_0;B_1;R;R
\]
that halts from every natural command-boundary work-store, possibly earlier
than its final command. Its first command is $\mathrm{AX}(0;0,0)$, with
opcode $2$, so its exact positive-continuation code is positive. The literal
initializer leaves next candidate $1$ and an empty residual. Ordinary
increment-by-one enumeration therefore reaches this one fixed successful
candidate unless it has already halted. The varying-parameter
$\mathrm{START}_Y$ family is unnecessary for this particular cutoff.

An opcode-independent version follows from an already established fact:
the credit proof gives $|W(n)|\le n+1$. The exact codec round trip, together
with $|\mathrm{STORE}_0|=10$, yields
\[
 \operatorname{code}(U)\ge |U|-1=21+|B_1|\ge21.
\]
Thus no new length induction is needed for the literal application. More
generally, for any codec satisfying
$\operatorname{decode}(\operatorname{pack}(c,r))=(c,r)$ and
$\operatorname{rest}(\operatorname{decode}(z))\le z$, one obtains
$r\le\operatorname{pack}(c,r)$, and induction on the positive-continuation
encoding proves $\operatorname{code}(v)\ge |v|-1$ for nonempty $v$.
Prefixing a universally halting word by $N$ arbitrary commands gives a
universally halting word with code at least $N$, provided every prefix
command terminates normally or halts. This is a reusable arbitrary-cutoff
lemma; it gives no hitting theorem for a skipping scheduler.

The latest phase-four \texttt{EnumPerm} already has a four-line specialized
code bound and a nine-line general completeness proof. The generic lemma is
therefore an abstraction opportunity, not a demonstrated local source-size
or compilation-time reduction. No sealed proof was changed and the proposed
Lean adapter has not been compiled.
